\sectionheading{5.48 অধ্যায়ের সারসংক্ষেপ ও পরীক্ষাভিত্তিক প্রস্তুতি}

এই অংশে পুরো অধ্যায়ের গুরুত্বপূর্ণ ধারণা, সংজ্ঞাভিত্তিক সংক্ষিপ্ত প্রশ্ন,
CQ-তে সম্ভাব্য প্রশ্নের ধরন, Programming Problem চিনে সমাধানের কৌশল,
Dry Run, Common Mistake এবং Last-minute Revision একসঙ্গে সাজানো হয়েছে।

\begin{rememberbox}

\begin{center}
\textbf{এই অংশটি নতুন Theory শেখার জন্য নয়; বরং পুরো Chapter দ্রুত Revision ও Exam Preparation-এর জন্য।}
\end{center}

\end{rememberbox}

%================================================
% 1
%================================================
\booksubheading{১. Chapter Overview — পুরো অধ্যায় এক নজরে}

\begin{revisionbox}{Programming Language থেকে Python পর্যন্ত সম্পর্ক}

\begin{center}

\textbf{Program}

\(\downarrow\)

\textbf{Programming Language}

\(\downarrow\)

Machine Language
\(\rightarrow\)
Assembly Language
\(\rightarrow\)
Mid Level
\(\rightarrow\)
High Level
\(\rightarrow\)
4GL

\medskip

\(\downarrow\)

\textbf{Translator Program}

Compiler
\(\rightarrow\)
Interpreter
\(\rightarrow\)
Assembler

\medskip

\(\downarrow\)

\textbf{Program Organization ও Design}

Program Design Model
\(\rightarrow\)
Modular Programming
\(\rightarrow\)
Program Development

\medskip

\(\downarrow\)

\textbf{Computational Thinking}

Problem Analysis
\(\rightarrow\)
IPO Model
\(\rightarrow\)
Dry Run / Trace Table

\medskip

\(\downarrow\)

\textbf{Algorithm}
\(\rightarrow\)
\textbf{Flowchart}

\medskip

\(\downarrow\)

\textbf{Python Programming}

Syntax
\(\rightarrow\)
Variable
\(\rightarrow\)
Data Type
\(\rightarrow\)
Input/Output
\(\rightarrow\)
Operator

\medskip

\(\downarrow\)

Conditional
\(\rightarrow\)
for Loop
\(\rightarrow\)
while Loop
\(\rightarrow\)
Loop Control
\(\rightarrow\)
Pattern

\medskip

\(\downarrow\)

List
\(\rightarrow\)
Tuple
\(\rightarrow\)
Set
\(\rightarrow\)
Dictionary

\medskip

\(\downarrow\)

Function
\(\rightarrow\)
Parameter
\(\rightarrow\)
Argument
\(\rightarrow\)
Return

\medskip

\(\downarrow\)

Function Contract
\(\rightarrow\)
Scope
\(\rightarrow\)
Local / Global
\(\rightarrow\)
File

\end{center}

\end{revisionbox}

\begin{rememberbox}

\begin{center}

\textbf{Problem}
\(\rightarrow\)
\textbf{Analysis}
\(\rightarrow\)
\textbf{Algorithm}
\(\rightarrow\)
\textbf{Flowchart}
\(\rightarrow\)
\textbf{Code}
\(\rightarrow\)
\textbf{Dry Run / Test}
\(\rightarrow\)
\textbf{Output}

\end{center}

\end{rememberbox}

%================================================
% 2
%================================================
\booksubheading{২. অধ্যায়ের মূল ধারণা — এক লাইনে}

\begin{longtable}{
|L{0.28\textwidth}
|L{0.64\textwidth}|
}

\twocoltablehead{বিষয়}{এক লাইনের সারকথা}
\endfirsthead

\twocoltablehead{বিষয়}{এক লাইনের সারকথা}
\endhead

Program
&
কম্পিউটারকে নির্দিষ্ট কাজ করানোর জন্য সাজানো নির্দেশনার সমষ্টি।
\\
\hline

Programming Language
&
Program লেখার জন্য ব্যবহৃত নিয়ম, Symbol ও Syntax-সমৃদ্ধ ভাষা।
\\
\hline

Machine Language
&
0 ও 1 ভিত্তিক কম্পিউটারের সরাসরি বোধগম্য ভাষা।
\\
\hline

Assembly Language
&
Mnemonic বা Symbolic Instruction-ভিত্তিক Low Level Language।
\\
\hline

High Level Language
&
মানুষের জন্য তুলনামূলক সহজ ও Problem-oriented Programming Language।
\\
\hline

4GL
&
কম Code ব্যবহার করে Application বা Database-ভিত্তিক কাজ সহজ করার উচ্চতর প্রজন্মের ভাষা।
\\
\hline

Translator Program
&
Source Program-কে অন্য ভাষায় রূপান্তরকারী Program।
\\
\hline

Compiler
&
পুরো Source Program একসঙ্গে অনুবাদ করে।
\\
\hline

Interpreter
&
সাধারণত Statement ধরে ধরে অনুবাদ ও Execute করে।
\\
\hline

Assembler
&
Assembly Language-কে Machine Language-এ রূপান্তর করে।
\\
\hline

Program Organization
&
Program-এর বিভিন্ন অংশকে পরিকল্পিতভাবে সাজানো।
\\
\hline

Modular Programming
&
বড় Program-কে ছোট ছোট Module-এ ভাগ করে তৈরি করার পদ্ধতি।
\\
\hline

Computational Thinking
&
সমস্যা বিশ্লেষণ করে যৌক্তিক ও বাস্তবায়নযোগ্য সমাধান তৈরির চিন্তাপদ্ধতি।
\\
\hline

IPO Model
&
Input
\(\rightarrow\)
Processing
\(\rightarrow\)
Output ভিত্তিক Problem Analysis।
\\
\hline

Dry Run
&
Program না চালিয়ে হাতে ধাপে ধাপে Execution যাচাই।
\\
\hline

Trace Table
&
Execution-এর বিভিন্ন ধাপে Variable-এর Value পরিবর্তন Table-এ দেখানো।
\\
\hline

Algorithm
&
সমস্যা সমাধানের সসীম, ধারাবাহিক ও সুনির্দিষ্ট ধাপ।
\\
\hline

Flowchart
&
Algorithm-এর Standard Symbol-ভিত্তিক চিত্ররূপ।
\\
\hline

Python
&
সহজ Syntax-সম্পন্ন High Level, General-purpose Programming Language।
\\
\hline

Syntax
&
Programming Language-এ Code লেখার নিয়ম।
\\
\hline

Variable
&
Data সংরক্ষণের জন্য ব্যবহৃত নাম।
\\
\hline

Data Type
&
কোনো Data কী ধরনের Value বহন করে তার শ্রেণি।
\\
\hline

Input
&
Program-এ Data গ্রহণ।
\\
\hline

Output
&
Program-এর Result প্রদর্শন।
\\
\hline

Operator
&
Operand-এর উপর Operation সম্পাদনকারী Symbol বা Keyword।
\\
\hline

Conditional Statement
&
Condition-এর ফল অনুযায়ী ভিন্ন কাজ সম্পাদনের Statement।
\\
\hline

Loop
&
একই কাজ প্রয়োজন অনুযায়ী বারবার সম্পাদনের পদ্ধতি।
\\
\hline

\py{break}
&
চলমান Loop সম্পূর্ণ বন্ধ করে।
\\
\hline

\py{continue}
&
বর্তমান Iteration বাদ দিয়ে পরবর্তী Iteration-এ যায়।
\\
\hline

\py{pass}
&
কোনো কাজ না করে Placeholder হিসেবে থাকে।
\\
\hline

List
&
Ordered ও Mutable Collection।
\\
\hline

Tuple
&
Ordered কিন্তু Immutable Collection।
\\
\hline

Set
&
Unique Element সংরক্ষণকারী Collection।
\\
\hline

Dictionary
&
Key--Value Pair-ভিত্তিক Collection।
\\
\hline

Function
&
নির্দিষ্ট কাজের জন্য reusable Code Block।
\\
\hline

Parameter
&
Function Definition-এ ব্যবহৃত Input Variable।
\\
\hline

Argument
&
Function Call-এর সময় পাঠানো Actual Value।
\\
\hline

Return Value
&
Function থেকে Caller-এর কাছে ফেরত পাঠানো Result।
\\
\hline

Function Contract
&
Function কী Input নেবে, কী করবে এবং কী Output দেবে তার নিয়ম।
\\
\hline

Scope
&
Program-এর কোন অংশে কোনো Variable বা Name ব্যবহারযোগ্য তার সীমা।
\\
\hline

Local Variable
&
সাধারণত Function-এর ভিতরে তৈরি ও সেই Function-এর Scope-এ ব্যবহৃত Variable।
\\
\hline

Global Variable
&
Function-এর বাইরে Module Level-এ তৈরি Variable।
\\
\hline

File
&
Program শেষ হওয়ার পরও Data সংরক্ষণের মাধ্যম।
\\
\hline

\end{longtable}

%================================================
% 3
%================================================
\booksubheading{৩. খুব সংক্ষিপ্ত প্রশ্ন — Definition / এক কথায় উত্তর}

এই প্রশ্নগুলো Definition, এক কথায় উত্তর, মৌখিক Revision এবং
জ্ঞানমূলক প্রশ্নের প্রস্তুতির জন্য উপযোগী।

\begin{longtable}{
|L{0.33\textwidth}
|L{0.59\textwidth}|
}

\twocoltablehead{প্রশ্ন}{সংক্ষিপ্ত উত্তর}
\endfirsthead

\twocoltablehead{প্রশ্ন}{সংক্ষিপ্ত উত্তর}
\endhead

১. Program কী?
&
কম্পিউটারকে নির্দিষ্ট কাজ করানোর জন্য সাজানো নির্দেশনার সমষ্টি।
\\
\hline

২. Programming Language কী?
&
Program লেখার জন্য ব্যবহৃত নিয়ম, Symbol ও Syntax-সমৃদ্ধ ভাষা।
\\
\hline

৩. Machine Language কী?
&
0 ও 1 দিয়ে গঠিত কম্পিউটারের সরাসরি বোধগম্য ভাষা।
\\
\hline

৪. Assembly Language কী?
&
Mnemonic Code-ভিত্তিক Low Level Programming Language।
\\
\hline

৫. Translator Program কী?
&
এক ভাষায় লেখা Program-কে অন্য ভাষায় রূপান্তরকারী Program।
\\
\hline

৬. Compiler কী?
&
পুরো Source Program একসঙ্গে অনুবাদকারী Translator।
\\
\hline

৭. Interpreter কী?
&
সাধারণত Statement ধরে ধরে অনুবাদ ও Execute করা Translator।
\\
\hline

৮. Assembler কী?
&
Assembly Language-কে Machine Language-এ রূপান্তরকারী Translator।
\\
\hline

৯. Modular Programming কী?
&
বড় Program-কে ছোট ছোট Module-এ ভাগ করে তৈরির পদ্ধতি।
\\
\hline

১০. Computational Thinking কী?
&
সমস্যাকে বিশ্লেষণ করে যৌক্তিক ও বাস্তবায়নযোগ্য ধাপে সমাধান তৈরির চিন্তাপদ্ধতি।
\\
\hline

১১. IPO-এর পূর্ণরূপ কী?
&
Input -- Processing -- Output।
\\
\hline

১২. Dry Run কী?
&
Program Execute না করে হাতে ধাপে ধাপে Logic যাচাই করার প্রক্রিয়া।
\\
\hline

১৩. Trace Table কী?
&
Execution-এর বিভিন্ন ধাপে Variable-এর Value পরিবর্তন Table-এ দেখানোর পদ্ধতি।
\\
\hline

১৪. Algorithm কী?
&
সমস্যা সমাধানের সসীম ও সুনির্দিষ্ট ধাপসমূহ।
\\
\hline

১৫. Flowchart কী?
&
Algorithm-এর Standard Symbol-ভিত্তিক চিত্ররূপ।
\\
\hline

১৬. Python কী?
&
একটি High Level, General-purpose Programming Language।
\\
\hline

১৭. Syntax কী?
&
Programming Language-এ সঠিকভাবে Code লেখার নিয়ম।
\\
\hline

১৮. Variable কী?
&
Data সংরক্ষণের জন্য ব্যবহৃত নাম।
\\
\hline

১৯. Data Type কী?
&
Data-এর ধরন বা প্রকৃতি নির্দেশকারী শ্রেণি।
\\
\hline

২০. Operator কী?
&
Operand-এর উপর Operation সম্পাদনকারী Symbol বা Keyword।
\\
\hline

২১. Conditional Statement কী?
&
Condition-এর ফল অনুযায়ী ভিন্ন Statement Execute করার পদ্ধতি।
\\
\hline

২২. Loop কী?
&
কোনো কাজ বারবার সম্পাদনের Programming Structure।
\\
\hline

২৩. \py{break} কী করে?
&
চলমান Loop সম্পূর্ণ বন্ধ করে।
\\
\hline

২৪. \py{continue} কী করে?
&
বর্তমান Iteration বাদ দিয়ে পরবর্তী Iteration শুরু করে।
\\
\hline

২৫. \py{pass} কী?
&
কোনো Operation না করা Placeholder Statement।
\\
\hline

২৬. List কী?
&
Ordered ও Mutable Collection।
\\
\hline

২৭. Tuple কী?
&
Ordered কিন্তু Immutable Collection।
\\
\hline

২৮. Set কী?
&
Unique Element সংরক্ষণকারী Collection।
\\
\hline

২৯. Dictionary কী?
&
Key--Value Pair আকারে Data সংরক্ষণকারী Collection।
\\
\hline

৩০. Function কী?
&
নির্দিষ্ট কাজ সম্পন্ন করার জন্য তৈরি reusable Code Block।
\\
\hline

৩১. Parameter কী?
&
Function Definition-এ Input গ্রহণকারী Variable।
\\
\hline

৩২. Argument কী?
&
Function Call-এর সময় Parameter-এর জন্য পাঠানো Actual Value।
\\
\hline

৩৩. Return Value কী?
&
Function থেকে Caller-এর কাছে ফেরত পাঠানো Result।
\\
\hline

৩৪. Function Contract কী?
&
Function-এর Input, Processing এবং Expected Output সম্পর্কিত নির্ধারিত নিয়ম।
\\
\hline

৩৫. Scope কী?
&
Program-এর যে অংশে কোনো Variable বা Name ব্যবহারযোগ্য সেই অঞ্চল।
\\
\hline

৩৬. Local Variable কী?
&
সাধারণত Function-এর ভিতরে তৈরি ও সেই Function-এর Scope-এ ব্যবহৃত Variable।
\\
\hline

৩৭. Global Variable কী?
&
Function-এর বাইরে Module Level-এ তৈরি Variable।
\\
\hline

৩৮. \py{global} Keyword কেন ব্যবহৃত হয়?
&
Function-এর ভিতরে Global Variable-এ Assignment করার জন্য।
\\
\hline

৩৯. File কী?
&
Persistentভাবে Data সংরক্ষণের মাধ্যম।
\\
\hline

৪০. Variable ও File-এর মূল পার্থক্য কী?
&
Variable সাধারণত Runtime Data রাখে; File Program শেষ হওয়ার পরও Data রাখতে পারে।
\\
\hline

\end{longtable}

%================================================
% 4
%================================================
\booksubheading{৪. “কেন / ব্যাখ্যা কর” ধরনের সংক্ষিপ্ত প্রশ্ন}

\begin{longtable}{
|L{0.38\textwidth}
|L{0.54\textwidth}|
}

\twocoltablehead{প্রশ্ন}{উত্তরের মূল ধারণা}
\endfirsthead

\twocoltablehead{প্রশ্ন}{উত্তরের মূল ধারণা}
\endhead

Compiler ও Interpreter-এর পার্থক্য কেন গুরুত্বপূর্ণ?
&
Compiler পুরো Program একসঙ্গে Translate করে; Interpreter সাধারণত Statement ধরে Translate/Execute করে। ফলে Translation ও Error handling-এর পদ্ধতি ভিন্ন।
\\
\hline

Algorithm আগে তৈরি করা হয় কেন?
&
Code লেখার আগে Problem-এর Logic পরিষ্কার হয় এবং ভুল কমে।
\\
\hline

Flowchart ব্যবহার করা হয় কেন?
&
Program Logic দৃশ্যমানভাবে বোঝা, বিশ্লেষণ এবং ব্যাখ্যা করা সহজ হয়।
\\
\hline

IPO Model কেন গুরুত্বপূর্ণ?
&
Problem-এর Input, Processing এবং Expected Output আলাদাভাবে শনাক্ত করতে সাহায্য করে।
\\
\hline

Dry Run কেন করা হয়?
&
Program Execute করার আগেই Logic ও Variable-এর পরিবর্তন যাচাই করা যায়।
\\
\hline

Python-এ Indentation গুরুত্বপূর্ণ কেন?
&
Indentation দ্বারা Python-এর Code Block নির্ধারিত হয়।
\\
\hline

\py{input()}-এর Result অনেক সময় Convert করতে হয় কেন?
&
\py{input()} সাধারণত String দেয়; Arithmetic করার জন্য প্রয়োজন অনুযায়ী \py{int()} বা \py{float()} দরকার হয়।
\\
\hline

\py{range(1, n+1)}-এ \py{n+1} কেন?
&
\py{range()}-এর Stop Value অন্তর্ভুক্ত হয় না; তাই \(n\) পর্যন্ত Value পেতে \py{n+1} দেওয়া হয়।
\\
\hline

Factorial-এর ক্ষেত্রে \py{fact = 1} নেওয়া হয় কেন?
&
Factorial-এ বারবার গুণ হয় এবং Multiplication-এর Neutral Initial Value হলো 1।
\\
\hline

Sum-এর ক্ষেত্রে \py{total = 0} নেওয়া হয় কেন?
&
Addition-এর Neutral Initial Value হলো 0।
\\
\hline

Function ব্যবহার করা ভালো কেন?
&
Code reuse, modularity, readability এবং maintenance সহজ হয়।
\\
\hline

\py{return} কেন গুরুত্বপূর্ণ?
&
Function-এর Result Caller-এর কাছে ফেরত দেয়; ফলে Result অন্য জায়গায় ব্যবহার করা যায়।
\\
\hline

Global Variable কম ব্যবহার করা ভালো কেন?
&
অপ্রয়োজনীয় Global Variable Program-এর বিভিন্ন অংশের মধ্যে Dependency বাড়াতে পারে।
\\
\hline

File কেন প্রয়োজন?
&
Program বন্ধ হওয়ার পরও Data সংরক্ষণ বা পরে আবার ব্যবহার করার প্রয়োজন হতে পারে।
\\
\hline

\end{longtable}

%================================================
% 5
%================================================
\booksubheading{৫. গুরুত্বপূর্ণ পার্থক্য — এক নজরে}

\begin{longtable}{
|L{0.23\textwidth}
|L{0.335\textwidth}
|L{0.335\textwidth}|
}

\threecoltablehead{তুলনা}{প্রথমটি}{দ্বিতীয়টি}
\endfirsthead

\threecoltablehead{তুলনা}{প্রথমটি}{দ্বিতীয়টি}
\endhead

Machine vs Assembly
&
Binary Code
&
Mnemonic / Symbolic Code
\\
\hline

Low Level vs High Level
&
Machine-oriented
&
User / Problem-oriented
\\
\hline

Compiler vs Interpreter
&
পুরো Program একসঙ্গে Translate
&
সাধারণত Statement ধরে Translate/Execute
\\
\hline

Algorithm vs Flowchart
&
লিখিত ধাপ
&
চিত্রভিত্তিক উপস্থাপন
\\
\hline

\py{=} vs \py{==}
&
Assignment
&
Comparison
\\
\hline

\py{/} vs \py{//}
&
সাধারণ Division
&
Floor Division
\\
\hline

\py{for} vs \py{while}
&
Range/Iterable-ভিত্তিক পুনরাবৃত্তিতে সুবিধাজনক
&
Condition-ভিত্তিক পুনরাবৃত্তিতে সুবিধাজনক
\\
\hline

\py{break} vs \py{continue}
&
Loop বন্ধ
&
বর্তমান Iteration Skip
\\
\hline

List vs Tuple
&
Mutable
&
Immutable
\\
\hline

List vs Set
&
Ordered, Duplicate থাকতে পারে
&
Unique Element
\\
\hline

List vs Dictionary
&
Element Collection
&
Key--Value Pair
\\
\hline

Parameter vs Argument
&
Function Definition-এ
&
Function Call-এ
\\
\hline

\py{print()} vs \py{return}
&
Output প্রদর্শন
&
Function থেকে Value ফেরত
\\
\hline

Local vs Global
&
Function-এর Local Scope
&
Module-level Scope
\\
\hline

Variable vs File
&
Runtime Data
&
Persistent Data
\\
\hline

\end{longtable}

%================================================
% 6
%================================================
\booksubheading{৬. CQ Question-এর ধরন বুঝে উত্তর দাও}

\begin{definitionbox}

CQ-এর প্রতিটি অংশে একই ধরনের কাজ চাওয়া হয় না। প্রশ্নের নির্দেশনা পড়ে
কোন Skill যাচাই করা হচ্ছে তা আগে বুঝতে হবে।

\end{definitionbox}

\begin{tabularx}{\linewidth}{
|C{0.10\linewidth}
|L{0.28\linewidth}
|Y|
}
\hline

\rowcolor{softBlue}
\textbf{\color{darkTeal}অংশ}
&
\textbf{\color{darkTeal}সাধারণ প্রকৃতি}
&
\textbf{\color{darkTeal}কীভাবে প্রস্তুতি নেবে}
\\
\hline

\textbf{ক}
&
জ্ঞানমূলক
&
সংজ্ঞা, পূর্ণরূপ, সরাসরি ধারণা ও এক কথায় উত্তর।
\\
\hline

\textbf{খ}
&
অনুধাবনমূলক
&
ব্যাখ্যা, কারণ, পার্থক্য, সুবিধা বা Concept বোঝানোর প্রস্তুতি।
\\
\hline

\textbf{গ}
&
প্রয়োগমূলক
&
উদ্দীপকের তথ্য ব্যবহার করে Algorithm, Flowchart, Code, Output বা প্রয়োগ।
\\
\hline

\textbf{ঘ}
&
উচ্চতর দক্ষতা
&
বিশ্লেষণ, তুলনা, পরিবর্তন, Logic, যুক্তি বা বিকল্প সমাধান।
\\
\hline

\end{tabularx}

\begin{rememberbox}

\textbf{\color{deepOrange}ভুল ধারণা:}

\begin{center}
“গ মানেই Program” অথবা “ঘ মানেই Program” — এটি ঠিক নয়।
\end{center}

প্রশ্নে কী চাওয়া হয়েছে সেটিই অনুসরণ করবে।

\end{rememberbox}

%================================================
% 7
%================================================
\booksubheading{৭. CQ-তে যেসব Programming Pattern অনুশীলন করবে}

\begin{tabularx}{\linewidth}{
|L{0.27\linewidth}
|Y|
}
\hline

\rowcolor{softBlue}
\textbf{\color{darkTeal}Problem Type}
&
\textbf{\color{darkTeal}গুরুত্বপূর্ণ Practice}
\\
\hline

Arithmetic / Formula
&
Sum, Average, Area, Perimeter, Temperature Conversion, Unit Conversion
\\
\hline

Conditional
&
Even/Odd, Positive/Negative, Largest/Smallest, Grade, Divisibility, Leap Year
\\
\hline

Loop
&
1 থেকে n, Odd/Even Number, Sum, Factorial, Repeated Output
\\
\hline

Series
&
\(1+2+\cdots+n\), Square Series, Odd/Even Series, Power Series
\\
\hline

Nested Loop
&
Number Pattern, Symbol Pattern
\\
\hline

List
&
Traversal, Sum, Average, Largest/Smallest, Search
\\
\hline

Function
&
Definition, Parameter, Argument, Return, Function Call
\\
\hline

Dry Run
&
Loop Output, Variable Update, Conditional Result
\\
\hline

Algorithm / Flowchart
&
Problem থেকে Algorithm, Algorithm থেকে Flowchart, Flowchart থেকে Program
\\
\hline

Code Analysis
&
Output নির্ণয়, Error শনাক্ত, Logic ব্যাখ্যা
\\
\hline

\end{tabularx}

%================================================
% 8
%================================================
\booksubheading{৮. Question-এর Keyword দেখেই কী করবে}

\begin{longtable}{
|L{0.30\textwidth}
|L{0.62\textwidth}|
}

\twocoltablehead{Question-এর শব্দ}{Answer-এর ধরন}
\endfirsthead

\twocoltablehead{Question-এর শব্দ}{Answer-এর ধরন}
\endhead

কী / কাকে বলে
&
সরাসরি Definition।
\\
\hline

পূর্ণরূপ লিখ
&
Full Form।
\\
\hline

ব্যাখ্যা কর
&
মূল বক্তব্য + সংক্ষিপ্ত ব্যাখ্যা।
\\
\hline

কেন
&
কারণ লিখবে।
\\
\hline

পার্থক্য লিখ
&
Table বা Point-wise Difference।
\\
\hline

তুলনা কর
&
দুটি Concept-এর বৈশিষ্ট্য পাশাপাশি লিখবে।
\\
\hline

সুবিধা লিখ
&
প্রাসঙ্গিক Advantage।
\\
\hline

নির্ণয় কর
&
Calculation বা Logic দেখাবে।
\\
\hline

Algorithm লেখ
&
Serial Step লিখবে।
\\
\hline

Flowchart আঁক
&
Standard Symbol ও Arrow ব্যবহার করবে।
\\
\hline

Program লেখ
&
Complete Logic + Correct Python Syntax।
\\
\hline

Output নির্ণয় কর
&
Dry Run / Trace করবে।
\\
\hline

বিশ্লেষণ কর
&
Result + Logic + কারণ।
\\
\hline

যৌক্তিকতা দেখাও
&
সিদ্ধান্ত + Supporting Reasoning।
\\
\hline

\end{longtable}

%================================================
% 9
%================================================
\booksubheading{৯. Problem দেখে কোন Concept ব্যবহার করবে?}

\begin{tabularx}{\linewidth}{
|L{0.40\linewidth}
|Y|
}
\hline

\rowcolor{softBlue}
\textbf{\color{darkTeal}Problem-এর Signal}
&
\textbf{\color{darkTeal}সম্ভাব্য Concept}
\\
\hline

“যদি”, “শর্ত”, “বড়/ছোট”
&
\py{if}, \py{elif}, \py{else}
\\
\hline

“বারবার”, “১ থেকে n”
&
Loop
\\
\hline

নির্দিষ্ট Range
&
\py{for} + \py{range()}
\\
\hline

Condition True থাকা পর্যন্ত
&
\py{while}
\\
\hline

Loop সঙ্গে সঙ্গে বন্ধ
&
\py{break}
\\
\hline

একটি Iteration বাদ
&
\py{continue}
\\
\hline

অনেক Data একসঙ্গে
&
List / Tuple / Set / Dictionary
\\
\hline

একই কাজ বারবার দরকার
&
Function
\\
\hline

Result অন্য জায়গায় ব্যবহার
&
\py{return}
\\
\hline

Program বন্ধের পরও Data দরকার
&
File
\\
\hline

Variable কোথায় ব্যবহারযোগ্য
&
Scope
\\
\hline

Value ধাপে ধাপে পরিবর্তিত হচ্ছে
&
Dry Run / Trace Table
\\
\hline

\end{tabularx}

%================================================
% 10
%================================================
\booksubheading{১০. Algorithm থেকে Flowchart থেকে Python Shortcut}

\begin{revisionbox}{Algorithm তৈরির সবচেয়ে সহজ কৌশল}

প্রথমে তিনটি প্রশ্নের উত্তর বের করো—

\begin{center}

\textbf{Input কী?}

\(\downarrow\)

\textbf{Processing / Condition কী?}

\(\downarrow\)

\textbf{Output কী?}

\end{center}

তারপর Algorithm-এর সাধারণ কাঠামো—

\begin{center}

\textbf{Start}
\(\rightarrow\)
\textbf{Input}
\(\rightarrow\)
\textbf{Processing / Decision}
\(\rightarrow\)
\textbf{Output}
\(\rightarrow\)
\textbf{End}

\end{center}

\end{revisionbox}

\begin{tabularx}{\linewidth}{
|L{0.31\linewidth}
|L{0.27\linewidth}
|Y|
}
\hline

\rowcolor{softBlue}
\textbf{\color{darkTeal}কাজ}
&
\textbf{\color{darkTeal}Flowchart}
&
\textbf{\color{darkTeal}Python}
\\
\hline

Program শুরু / শেষ
&
Oval
&
Program-এর শুরু / শেষ
\\
\hline

Input / Output
&
Parallelogram
&
\py{input()}, \py{print()}
\\
\hline

Processing
&
Rectangle
&
Assignment / Arithmetic
\\
\hline

Decision
&
Diamond
&
\py{if}, \py{elif}, \py{else}
\\
\hline

Repetition
&
Looping Flow
&
\py{for}, \py{while}
\\
\hline

Loop বন্ধ
&
Flow Exit
&
\py{break}
\\
\hline

Iteration Skip
&
Next Iteration
&
\py{continue}
\\
\hline

\end{tabularx}

%================================================
% 11
%================================================
\booksubheading{১১. Dry Run / Output Question সমাধানের নিয়ম}

\begin{tabularx}{\linewidth}{
|>{\RaggedRight\arraybackslash}X
|>{\RaggedRight\arraybackslash}X|
}
\hline

\rowcolor{darkTeal}
\multicolumn{2}{|l|}{
\color{white}\bfseries উদাহরণ: Loop-এর Output Dry Run
}
\\
\hline

\rowcolor{softGreen}
\textbf{\color{darkTeal}Python Code}
&
\textbf{\color{darkTeal}Trace Table}
\\
\hline

\begin{minipage}[t]{\linewidth}
\vspace{0pt}

\begin{english}
\ttfamily
total = 0\\[0.5mm]
for i in range(1, 5):\\
\hspace*{4mm}total = total + i\\[0.5mm]
print(total)
\end{english}

\end{minipage}
&
\begin{minipage}[t]{\linewidth}
\vspace{0pt}

\centering

\begin{tabular}{|c|c|}
\hline
\textbf{i} & \textbf{total} \\
\hline
শুরু & 0 \\
\hline
1 & 1 \\
\hline
2 & 3 \\
\hline
3 & 6 \\
\hline
4 & 10 \\
\hline
\end{tabular}

\medskip

\textbf{Final Output = 10}

\end{minipage}
\\
\hline

\end{tabularx}

\begin{rememberbox}

\textbf{Dry Run Shortcut:}

\begin{center}

Initial Value
\(\rightarrow\)
Loop Range
\(\rightarrow\)
প্রতি Iteration-এর Update
\(\rightarrow\)
Condition
\(\rightarrow\)
Final Output

\end{center}

\end{rememberbox}

%================================================
% 12
%================================================
\booksubheading{১২. Code দেখে Output বের করার কৌশল}

\begin{enumerate}

\item Variable-এর Initial Value কত দেখো।

\item Input থাকলে কোন Value নেওয়া হয়েছে তা নির্ধারণ করো।

\item Condition True না False তা পরীক্ষা করো।

\item Loop কয়বার চলবে তা বের করো।

\item প্রতিবার কোন Variable পরিবর্তিত হচ্ছে লিখে রাখো।

\item \py{break} বা \py{continue} আছে কি না দেখো।

\item শেষে কোন Statement Print হচ্ছে তা নির্ধারণ করো।

\end{enumerate}

%================================================
% 13
%================================================
\booksubheading{১৩. Common Exam Trap — ভুল করো না}

\begin{trapbox}

\textbf{\color{deepOrange}১. \py{=} এবং \py{==}}

\py{=}
\(\rightarrow\)
Assignment

\py{==}
\(\rightarrow\)
Comparison

\medskip

\textbf{\color{deepOrange}২. \py{input()}}

\py{input()} সাধারণত String Return করে।

\begin{english}
\ttfamily
age = int(input("Age: "))
\end{english}

\medskip

\textbf{\color{deepOrange}৩. \py{range()}}

\py{range(1, 5)} দিলে পাওয়া যায়—

\py{1, 2, 3, 4}

অর্থাৎ Stop Value 5 অন্তর্ভুক্ত হয় না।

\medskip

\textbf{\color{deepOrange}৪. Sum-এর Initial Value}

\begin{english}
\ttfamily
total = 0
\end{english}

\medskip

\textbf{\color{deepOrange}৫. Factorial-এর Initial Value}

\begin{english}
\ttfamily
fact = 1
\end{english}

\medskip

\textbf{\color{deepOrange}৬. while Loop}

Update ভুলে গেলে Infinite Loop হতে পারে।

\medskip

\textbf{\color{deepOrange}৭. break বনাম continue}

\py{break}
\(\rightarrow\)
পুরো Loop বন্ধ।

\py{continue}
\(\rightarrow\)
শুধু বর্তমান Iteration বাদ।

\medskip

\textbf{\color{deepOrange}৮. List Index}

প্রথম Element-এর Index হলো \py{0}।

\medskip

\textbf{\color{deepOrange}৯. print() বনাম return}

\py{print()}
\(\rightarrow\)
Output দেখায়।

\py{return}
\(\rightarrow\)
Function থেকে Value ফেরত দেয়।

\medskip

\textbf{\color{deepOrange}১০. Local বনাম Global}

Function-এর Local Variable সাধারণত Function-এর বাইরে সরাসরি ব্যবহার করা যায় না।

\medskip

\textbf{\color{deepOrange}১১. Variable বনাম File}

Variable Runtime Data-এর জন্য; File Persistent Data সংরক্ষণে ব্যবহৃত হয়।

\medskip

\textbf{\color{deepOrange}১২. Power Operator}

Python-এ Power-এর জন্য \py{**} ব্যবহার হয়।

\end{trapbox}

%================================================
% 14
%================================================
\booksubheading{১৪. ৩০ সেকেন্ডে Chapter Revision}

\begin{revisionbox}{একবার চোখ বুলিয়ে নাও}

\begin{center}

Program
\(\rightarrow\)
Programming Language
\(\rightarrow\)
Machine / Assembly / High Level

\medskip

\(\downarrow\)

Translator
\(\rightarrow\)
Compiler / Interpreter / Assembler

\medskip

\(\downarrow\)

Program Design
\(\rightarrow\)
Computational Thinking
\(\rightarrow\)
IPO
\(\rightarrow\)
Dry Run

\medskip

\(\downarrow\)

Algorithm
\(\rightarrow\)
Flowchart

\medskip

\(\downarrow\)

Python Syntax
\(\rightarrow\)
Variable
\(\rightarrow\)
Data Type
\(\rightarrow\)
Input / Output
\(\rightarrow\)
Operator

\medskip

\(\downarrow\)

Conditional
\(\rightarrow\)
Loop
\(\rightarrow\)
Loop Control
\(\rightarrow\)
Pattern

\medskip

\(\downarrow\)

List
\(\rightarrow\)
Tuple
\(\rightarrow\)
Set
\(\rightarrow\)
Dictionary

\medskip

\(\downarrow\)

Function
\(\rightarrow\)
Parameter
\(\rightarrow\)
Argument
\(\rightarrow\)
Return

\medskip

\(\downarrow\)

Function Contract
\(\rightarrow\)
Scope
\(\rightarrow\)
Local / Global
\(\rightarrow\)
File

\end{center}

\end{revisionbox}

%================================================
% 15
%================================================
\booksubheading{১৫. Last-minute Revision Checklist}

\begin{tabularx}{\linewidth}{
|C{0.08\linewidth}
|Y|
}
\hline

\rowcolor{softBlue}
\textbf{\color{darkTeal}\(\checkmark\)}
&
\textbf{\color{darkTeal}পরীক্ষার আগে নিশ্চিত হও}
\\
\hline

\(\square\)
&
Program ও Programming Language-এর সংজ্ঞা বলতে পারি।
\\
\hline

\(\square\)
&
Machine, Assembly ও High Level Language-এর পার্থক্য জানি।
\\
\hline

\(\square\)
&
Compiler, Interpreter ও Assembler আলাদা করতে পারি।
\\
\hline

\(\square\)
&
IPO Model বুঝি।
\\
\hline

\(\square\)
&
Algorithm লিখতে পারি।
\\
\hline

\(\square\)
&
Flowchart-এর Standard Symbol চিনি।
\\
\hline

\(\square\)
&
Dry Run ও Trace Table করতে পারি।
\\
\hline

\(\square\)
&
Python Syntax ও Indentation বুঝি।
\\
\hline

\(\square\)
&
Variable ও Data Type বুঝি।
\\
\hline

\(\square\)
&
\py{input()} ও \py{print()} ব্যবহার করতে পারি।
\\
\hline

\(\square\)
&
Operator চিনতে পারি।
\\
\hline

\(\square\)
&
\py{if-elif-else} লিখতে পারি।
\\
\hline

\(\square\)
&
\py{for} ও \py{while} Loop লিখতে পারি।
\\
\hline

\(\square\)
&
\py{break}, \py{continue}, \py{pass} বুঝি।
\\
\hline

\(\square\)
&
List, Tuple, Set ও Dictionary আলাদা করতে পারি।
\\
\hline

\(\square\)
&
Function Define ও Call করতে পারি।
\\
\hline

\(\square\)
&
Parameter ও Argument-এর পার্থক্য জানি।
\\
\hline

\(\square\)
&
\py{return} কী করে বুঝি।
\\
\hline

\(\square\)
&
Local ও Global Variable বুঝি।
\\
\hline

\(\square\)
&
Function Contract ও Scope বলতে পারি।
\\
\hline

\(\square\)
&
File কেন প্রয়োজন বুঝি।
\\
\hline

\(\square\)
&
অন্তত ১০--১৫টি গুরুত্বপূর্ণ Program নিজে লিখতে পারি।
\\
\hline

\end{tabularx}

%================================================
% 16
%================================================
\booksubheading{১৬. Mini Self-Test}

\begin{revisionbox}{নিজেকে যাচাই করো}

\begin{enumerate}

\item Algorithm কী?

\item IPO-এর পূর্ণরূপ কী?

\item Compiler ও Interpreter-এর একটি মৌলিক পার্থক্য লেখ।

\item Python-এ \py{input()} সাধারণত কোন Data Type Return করে?

\item \py{break} ও \py{continue}-এর কাজ কী?

\item নিচের Code-এর Output কত?

\begin{english}
\ttfamily
x = 5\\[0.5mm]
if x > 3:\\
\hspace*{4mm}print("Yes")\\
else:\\
\hspace*{4mm}print("No")
\end{english}

\item নিচের Code-এর Output নির্ণয় কর।

\begin{english}
\ttfamily
total = 0\\[0.5mm]
for i in range(1, 4):\\
\hspace*{4mm}total += i\\[0.5mm]
print(total)
\end{english}

\item ১ থেকে \(n\) পর্যন্ত Even Number প্রদর্শনের Algorithm লেখ।

\item তিনটি সংখ্যার মধ্যে বৃহত্তম নির্ণয়ের একটি Decision Condition লেখ।

\item Function-এর Parameter ও Argument-এর মধ্যে পার্থক্য কী?

\end{enumerate}

\end{revisionbox}

%================================================
% 17
%================================================
\booksubheading{১৭. Mini Self-Test Answer Key}

\begin{longtable}{
|C{0.08\textwidth}
|L{0.84\textwidth}|
}

\twocoltablehead{নং}{উত্তর}
\endfirsthead

\twocoltablehead{নং}{উত্তর}
\endhead

১
&
সমস্যা সমাধানের সসীম ও সুনির্দিষ্ট ধাপসমূহ।
\\
\hline

২
&
Input -- Processing -- Output।
\\
\hline

৩
&
Compiler পুরো Program একসঙ্গে Translate করে; Interpreter সাধারণত Statement ধরে Translate/Execute করে।
\\
\hline

৪
&
\py{str}
\\
\hline

৫
&
\py{break} Loop বন্ধ করে; \py{continue} বর্তমান Iteration Skip করে।
\\
\hline

৬
&
\py{Yes}
\\
\hline

৭
&
\py{6}
\\
\hline

৮
&
Start
\(\rightarrow\)
\(n\) Input
\(\rightarrow\)
\(i=2\)
\(\rightarrow\)
\(i\le n\) হলে Print
\(\rightarrow\)
\(i=i+2\)
\(\rightarrow\)
Repeat
\(\rightarrow\)
End।
\\
\hline

৯
&
উদাহরণ:
\py{a >= b and a >= c}
\\
\hline

১০
&
Parameter Function Definition-এ থাকে; Argument Function Call-এর সময় পাঠানো Actual Value।
\\
\hline

\end{longtable}

%================================================
% Final Revision
%================================================
\begin{rememberbox}

\textbf{\color{deepOrange}শেষ মুহূর্তের প্রস্তুতির সূত্র}

\begin{center}

\textbf{সংজ্ঞা জানো}

\(\downarrow\)

\textbf{পার্থক্য বোঝো}

\(\downarrow\)

\textbf{Problem-এর Logic ধরো}

\(\downarrow\)

\textbf{Algorithm লেখো}

\(\downarrow\)

\textbf{Flowchart আঁকো}

\(\downarrow\)

\textbf{Python Code লেখো}

\(\downarrow\)

\textbf{Dry Run করে Output যাচাই করো}

\end{center}

\end{rememberbox}

\small

